\documentclass[11pt,a4paper]{article}

\usepackage[utf8]{inputenc}
\usepackage{amsmath,amssymb,amsfonts}
\usepackage{geometry}
\geometry{left=2.5cm, right=2.5cm, top=2.8cm, bottom=2.8cm}
\usepackage{xcolor}
\usepackage{booktabs}
\usepackage{fancyhdr}
\usepackage{listings}
\usepackage{hyperref}
\usepackage{enumitem}

% Paleta institucional UAN
\definecolor{uanblue}{RGB}{0, 45, 98}
\definecolor{uancyan}{RGB}{0, 140, 200}
\definecolor{accentgreen}{RGB}{0, 160, 120}
\definecolor{codebg}{RGB}{246, 248, 252}
\definecolor{bordercolor}{RGB}{205, 215, 230}
\definecolor{darkgray}{RGB}{60, 64, 70}

\hypersetup{
    colorlinks=true,
    linkcolor=uanblue,
    citecolor=uancyan,
    urlcolor=uancyan,
    pdftitle={Aurora By Juan Esteban Vera - Manual de Usuario e Interpretacion},
    pdfauthor={Juan Esteban Vera Rodriguez}
}

\pagestyle{fancy}
\fancyhf{}
\fancyhead[L]{\small \textcolor{uanblue}{\textbf{Universidad Antonio Nariño}} · Sede Ibagué · Ingeniería Electrónica}
\fancyhead[R]{\small \textcolor{gray}{Aurora By Juan Esteban Vera · Manual de Usuario}}
\fancyfoot[L]{\small \textcolor{gray}{Anexo Técnico 03 — Manual de Operación e Interpretación}}
\fancyfoot[R]{\small \thepage}
\renewcommand{\headrulewidth}{0.6pt}
\renewcommand{\footrulewidth}{0.4pt}

\lstset{
    backgroundcolor=\color{codebg},
    basicstyle=\ttfamily\footnotesize,
    breaklines=true,
    captionpos=b,
    frame=single,
    rulecolor=\color{bordercolor},
    keywordstyle=\color{uanblue}\bfseries,
    commentstyle=\color{gray}\itshape,
    stringstyle=\color{accentgreen},
    showstringspaces=false
}

\begin{document}

% =========================================================================
% PORTADA INSTITUCIONAL
% =========================================================================
\begin{titlepage}
    \centering
    \vspace*{0.5cm}
    
    {\scshape\LARGE \textbf{UNIVERSIDAD ANTONIO NARIÑO}}\\[0.3cm]
    {\scshape\large FACULTAD DE INGENIERÍA · PROGRAMA DE INGENIERÍA ELECTRÓNICA}\\[0.15cm]
    {\small SEDE IBAGUÉ — CAPITAL MUSICAL DE COLOMBIA}\\[1.5cm]
    
    \rule{\linewidth}{1.8pt}\\[0.4cm]
    {\Huge \bfseries \textcolor{uanblue}{MANUAL OFICIAL DE USUARIO}}\\[0.3cm]
    {\Large \bfseries GUÍA COMPLETA DE OPERACIÓN, INTERPRETACIÓN Y CONTROL DSP}\\[0.25cm]
    {\large \itshape Aurora By Juan Esteban Vera — Yamaha MONTAGE DSP Workstation}\\[0.4cm]
    \rule{\linewidth}{1.8pt}\\[2.0cm]
    
    \begin{minipage}{0.85\textwidth}
        \large
        \textbf{Autor y Desarrollador:}\\[0.15cm]
        \textbf{Juan Esteban Vera Rodríguez}\\
        Código Estudiantil / Investigador Principal\\
        Programa de Ingeniería Electrónica · Sede Ibagué\\[0.8cm]
        
        \textbf{Proyecto Institucional:}\\
        Aurora By Juan Esteban Vera — DSP Workstation 2.1.0\\
        Dedicado a la Universidad Antonio Nariño (UAN)\\[0.8cm]
        
        \textbf{Clasificación Documental:}\\
        Anexo Técnico ANX-USR-003 · Distribución Abierta y Académica
    \end{minipage}
    
    \vfill
    {\large Ibagué, Tolima, Colombia · 2026}
\end{titlepage}

\newpage
\tableofcontents
\newpage

% =========================================================================
\section{Introducción al Instrumento Aurora}
% =========================================================================

\textbf{Aurora By Juan Esteban Vera} es una estación de trabajo de síntesis y procesamiento digital de señales (DSP) en tiempo real, concebida para brindar la máxima fidelidad acústica, calidez analógica y versatilidad tímbrica tanto en el estudio de producción musical como en presentaciones en vivo.

El instrumento fusiona la calidez del \textbf{modelado analógico sustractivo} clásico (inspirado en el legendario Yamaha CS-80) con la riqueza espectral de la \textbf{síntesis por Modulación de Frecuencia} (Yamaha DX7 / FM-X), coronado por el célebre \textbf{Motion Control System con Super Knob} del sintetizador insignia \textbf{Yamaha MONTAGE}.

\subsection{Características Principales}
\begin{itemize}[leftmargin=1.5cm]
    \item \textbf{Polifonía Dinámica de 8 Voces:} Cada voz dispone de 2 osciladores con formas de onda sinusoidales, triangulares, diente de sierra y cuadradas con anti-aliasing PolyBLEP, además de generador de ruido blanco estocástico.
    \item \textbf{Motor de Frecuencia Modulada (FM-X):} Modulación de fase cruzada en tiempo real entre el Oscilador 2 (modulador) y el Oscilador 1 (portadora).
    \item \textbf{Filtro Resonante Analógico Virtual TPT SVF:} Filtro de variable de estado de 2 polos basado en la Transformada Bilineal con preservación de topología (Topology-Preserving Transform), asegurando estabilidad absoluta sin desbordamiento numérico.
    \item \textbf{Motion Control Super Knob:} Dial maestro central que interpola y modula simultáneamente múltiples parámetros (corte de filtro, resonancia, índice FM, drive y mezcla de efectos).
    \item \textbf{Secuenciador por Pasos de 16 Eventos:} Sincronizado por reloj interno con memoria de patrones y persistencia en formato JSON.
    \item \textbf{Procesador Vocal con Corrección Tonal (Autotune):} Módulo de captura por micrófono con análisis FFT y pitch-shifting en tiempo real hacia escalas musicales seleccionables.
    \item \textbf{Osciloscopio y Vúmetros Estéreo en Vivo:} Visualización continua de la forma de onda de salida a 30 FPS y medición de nivel de pico con amortiguamiento balístico.
\end{itemize}

% =========================================================================
\section{Requisitos del Sistema e Instalación}
% =========================================================================

\subsection{Requisitos de Hardware y Software}
\begin{itemize}[leftmargin=1.5cm]
    \item \textbf{Sistema Operativo:} Windows 10 / 11 (64 bits), macOS 11+ o GNU/Linux (Ubuntu 20.04+).
    \item \textbf{Procesador (CPU):} Intel Core i3 / AMD Ryzen 3 o superior (optimizado para SIMD AVX2/NEON vía NumPy).
    \item \textbf{Memoria RAM:} Mínimo 2 GB (el software en ejecución consume menos de $45\text{ MB}$).
    \item \textbf{Dispositivo de Sonido:} Dispositivo de audio estéreo estándar compatible con WASAPI, ASIO o CoreAudio para baja latencia.
    \item \textbf{Entorno de Ejecución:} Python 3.10 o superior con dependencias oficiales instaladas.
\end{itemize}

\subsection{Puesta en Marcha Rápida}
Para poner en marcha la estación de trabajo, siga estos sencillos pasos:

\begin{enumerate}[leftmargin=1.5cm]
    \item Abra una terminal de comandos (PowerShell, Command Prompt o Bash) en el directorio del proyecto.
    \item Verifique e instale las dependencias del sistema mediante:
\begin{lstlisting}[language=bash]
python -m pip install -r requirements.txt
\end{lstlisting}
    \item Ejecute la aplicación principal:
\begin{lstlisting}[language=bash]
python main.py
\end{lstlisting}
    \item La interfaz gráfica se desplegará instantáneamente, inicializando el motor de audio a $44.100\text{ Hz}$ con un bloque de 256 muestras y cargando el preset de bienvenida \textit{Yamaha DX7 E-Piano}.
\end{enumerate}

% =========================================================================
\section{Distribución del Teclado Extendido (25 Teclas)}
% =========================================================================

Aurora incorpora un teclado interactivo en pantalla de \textbf{25 teclas continuas} (15 blancas correspondientes a notas diatónicas y 10 negras para alteraciones cromáticas), abarcando dos octavas completas simultáneas:

\begin{center}
{\footnotesize
\setlength{\tabcolsep}{3.5pt}
\begin{tabular}{cccccccccccccccc}
\toprule
\textbf{Negras} & & C\#3 & D\#3 & & F\#3 & G\#3 & A\#3 & & C\#4 & D\#4 & & F\#4 & G\#4 & A\#4 & \\
\textbf{Blancas} & C3 & D3 & E3 & F3 & G3 & A3 & B3 & C4 & D4 & E4 & F4 & G4 & A4 & B4 & C5 \\
\bottomrule
\end{tabular}
}
\end{center}

\subsection{Interpretación con Ratón y Pantalla Táctil}
\begin{itemize}[leftmargin=1.5cm]
    \item \textbf{Clic Individual:} Dispara la nota instantáneamente con retroalimentación visual de color aguamarina.
    \item \textbf{Arrastre Glissando:} Manteniendo presionado el botón izquierdo mientras se arrastra el cursor por las teclas, se ejecutan barridos melódicos fluidos sin pérdidas de disparo.
\end{itemize}

\subsection{Mapeo del Teclado de la Computadora}
Para una ejecución musical cómoda y polifónica con ambas manos:

\begin{enumerate}[leftmargin=1.5cm]
    \item \textbf{Octava Inferior (Grave / Registro C3 a B3):}
    \begin{itemize}
        \item Notas Blancas: \texttt{Z} (Do3), \texttt{X} (Re3), \texttt{C} (Mi3), \texttt{V} (Fa3), \texttt{B} (Sol3), \texttt{N} (La3), \texttt{M} (Si3).
        \item Notas Negras: \texttt{S} (Do\#3), \texttt{D} (Re\#3), \texttt{G} (Fa\#3), \texttt{H} (Sol\#3), \texttt{J} (La\#3).
    \end{itemize}
    \item \textbf{Octava Superior (Aguda / Registro C4 a C5):}
    \begin{itemize}
        \item Notas Blancas: \texttt{Q} (Do4), \texttt{W} (Re4), \texttt{E} (Mi4), \texttt{R} (Fa4), \texttt{T} (Sol4), \texttt{Y} (La4), \texttt{U} (Si4), \texttt{I} (Do5).
        \item Notas Negras: \texttt{2} (Do\#4), \texttt{3} (Re\#4), \texttt{5} (Fa\#4), \texttt{6} (Sol\#4), \texttt{7} (La\#4).
    \end{itemize}
    \item \textbf{Mapeo Clásico Fila Central (C4 a C5):}
    \begin{itemize}
        \item \texttt{A}, \texttt{W}, \texttt{S}, \texttt{E}, \texttt{D}, \texttt{F}, \texttt{T}, \texttt{G}, \texttt{Y}, \texttt{H}, \texttt{U}, \texttt{J}, \texttt{K}.
    \end{itemize}
\end{enumerate}

\subsection{Control de Transposición de Octavas}
Los botones \textbf{\texttt{[ < OCT - ]}} y \textbf{\texttt{[ OCT + > ]}} permiten desplazar el rango sonoro completo en pasos de una octava:
\begin{itemize}[leftmargin=1.5cm]
    \item Transposición hacia abajo: Permite alcanzar registros de sub-bajos profundos hasta \textbf{C1} ($\approx 32.7\text{ Hz}$).
    \item Transposición hacia arriba: Permite alcanzar registros solistas brillantes de campana y flauta hasta \textbf{C7}.
\end{itemize}

\subsection{Botón de Pánico y Silenciado}
En caso de notas sostenidas indeseadas, presione la tecla \textbf{\texttt{Escape (Esc)}} o haga clic sobre \textbf{\texttt{[ SILENCIAR TODO ]}} en la cabecera para liberar de inmediato todas las voces activas del buffer.

% =========================================================================
\section{Guía Detallada de Módulos y Parámetros}
% =========================================================================

\subsection{Pestaña 1: Sintetizador \& FM-X (Yamaha DX7)}
\begin{itemize}[leftmargin=1.5cm]
    \item \textbf{Formas de Onda (OSC 1 y OSC 2):}
    \begin{itemize}
        \item \textit{Sine:} Seno matemático puro, fundamental para bajos redondos y bases de FM.
        \item \textit{Triangle:} Onda triangular suave rica en armónicos impares que decaen con $1/n^2$.
        \item \textit{Saw:} Diente de sierra brillante procesada con corrección PolyBLEP para eliminar aliasing.
        \item \textit{Square:} Onda cuadrada con armónicos impares fuertes, idónea para vientos y clarinetes.
        \item \textit{Noise:} Ruido estocástico con distribución uniforme, ideal para efectos de viento o percusión.
    \end{itemize}
    \item \textbf{Mezcla OSC 2 (\texttt{osc2\_mix}):} Rango $0.0$ a $1.0$. Controla la ponderación entre ambos osciladores.
    \item \textbf{Desafinación en Cents (\texttt{detune\_cents}):} Rango $\pm 40\text{ cents}$. Introduce un batimiento orgánico de coro analógico (\textit{chorus}).
    \item \textbf{Semitonos OSC 2 (\texttt{osc2\_semitones}):} Desplaza el segundo oscilador en intervalos de $-24$ a $+24$ semitonos (por ejemplo, $+7$ para quinta justa, $-12$ para sub-oscilador).
    \item \textbf{Profundidad FM (\texttt{fm\_depth}):} Modulación de fase de alta velocidad del OSC 1 gobernada por la salida del OSC 2. A valores bajos genera brillo vítreo; a valores altos produce texturas agresivas y campaniformes.
\end{itemize}

\subsection{Filtro Resonante VCF (Yamaha CS-80)}
\begin{itemize}[leftmargin=1.5cm]
    \item \textbf{Frecuencia de Corte (\texttt{cutoff}):} Rango $50\text{ Hz}$ a $14.000\text{ Hz}$. Atenúa armónicos por encima del umbral a una pendiente de $-12\text{ dB/octava}$.
    \item \textbf{Resonancia Q (\texttt{resonance}):} Rango $0.0$ a $0.95$. Genera un pico de realimentación acústica en la frecuencia de corte. El algoritmo TPT SVF satura suavemente el estado evitando cualquier inestabilidad matemática.
    \item \textbf{Vibrato LFO (\texttt{vibrato\_rate} y \texttt{vibrato\_depth}):} Modula la frecuencia de las voces entre $0.1\text{ Hz}$ y $15\text{ Hz}$ para dotar al sonido de calidez expresiva.
\end{itemize}

\subsection{Generador de Envolvente ADSR de Amplitud}
Permite esculpir la dinámica acústica temporal de cada pulsación:
\begin{itemize}[leftmargin=1.5cm]
    \item \textbf{Attack (Ataque):} Tiempo para alcanzar la amplitud máxima ($0.001\text{ s}$ a $2.5\text{ s}$).
    \item \textbf{Decay (Decaimiento):} Tiempo de caída desde el pico hasta el nivel sostenido ($0.01\text{ s}$ a $2.5\text{ s}$).
    \item \textbf{Sustain (Sostenimiento):} Amplitud normalizada de reposo mientras la tecla se mantenga pulsada ($0.0$ a $1.0$).
    \item \textbf{Release (Liberación):} Tiempo de extinción acústica al soltar la tecla ($0.01\text{ s}$ a $3.5\text{ s}$).
\end{itemize}

\subsection{Pestaña 2: Motion Control Super Knob \& Efectos MONTAGE}
\begin{itemize}[leftmargin=1.5cm]
    \item \textbf{El Super Knob Central:} Dial de modulación gestual macro que permite controlar la transición tímbrica en vivo con la rueda del ratón o arrastre vertical. A medida que sube:
    \begin{enumerate}
        \item Abre la frecuencia de corte del filtro progresivamente.
        \item Eleva la resonancia para mayor presencia armónica.
        \item Incrementa el índice de modulación FM.
        \item Aplica saturación cálida no lineal (\textit{Drive}).
        \item Expande la espacialidad del eco y la reverberación.
    \end{enumerate}
    \item \textbf{Analog Drive ($\tanh$):} Emula el comportamiento térmico de los bulbos o transistores saturados, enriqueciendo la señal con armónicos pares e impares.
    \item \textbf{Analog Stereo Delay:} Retardo estéreo con rebote ping-pong (desfase temporal de $15\%$ en el canal derecho) y realimentación protegida.
    \item \textbf{Reverberación Schroeder:} Red multi-tap con retardos primos que crea una sensación tridimensional de recinto acústico natural.
\end{itemize}

\subsection{Pestaña 3: Secuenciador por Pasos de 16 Eventos}
Permite la creación de arpegios, líneas de bajo y patrones rítmicos:
\begin{enumerate}[leftmargin=1.5cm]
    \item Ajuste el tempo de ejecución mediante el deslizador \textbf{BPM} ($50$ a $200\text{ BPM}$).
    \item Cada uno de los 16 botones representa una semicorchea. Al hacer clic se alterna secuencialmente la nota musical deseada o un silencio (\texttt{OFF}).
    \item Presione \textbf{\texttt{[ PLAY REPRODUCIR ]}} para iniciar la reproducción temporizada en hilo secundario.
    \item Utilice \textbf{\texttt{[ GUARDAR SECUENCIA ]}} para exportar el patrón junto con los ajustes de síntesis a un archivo \texttt{.aurora-song.json}, y \textbf{\texttt{[ ABRIR SECUENCIA ]}} para recuperarlo en cualquier momento.
\end{enumerate}

\subsection{Pestaña 4: Procesador Vocal \& Afinador Electrónico}
Herramienta de procesamiento de voz en vivo:
\begin{enumerate}[leftmargin=1.5cm]
    \item \textbf{Recomendación imprescindible:} Utilice audífonos para evitar el acople acústico (\textit{feedback}) con los altavoces.
    \item Marque la casilla \textbf{``Activar monitorización y afinación de micrófono''}.
    \item Seleccione la escala musical de referencia: \textit{C mayor}, \textit{C menor}, \textit{A menor} o \textit{C cromática}.
    \item Hable o cante hacia el micrófono: el visor indicará la frecuencia fundamental detectada y ajustará automáticamente el tono a la nota más próxima dentro de la tonalidad elegida.
    \item Gradúe el deslizador \textbf{Intensidad de Corrección} según el efecto deseado ($0.5$ para entonación natural, $1.0$ para el clásico efecto robótico de la música moderna).
\end{enumerate}

% =========================================================================
\section{Catálogo de Presets de Fábrica (Yamaha MONTAGE Edition)}
% =========================================================================

Aurora incluye 7 programas de fábrica cuidadosamente calibrados:
\begin{enumerate}[leftmargin=1.5cm]
    \item \textbf{Yamaha DX7 E-Piano:} Clásico piano eléctrico de síntesis FM cristalino y dinámico.
    \item \textbf{CS-80 Poly Brass:} Conjunto de metales analógicos majestuoso con filtro TPT cálido y amplio tiempo de ataque.
    \item \textbf{Montage Super Lead:} Solista cortante y agresivo con FM profunda y saturación analógica.
    \item \textbf{Ambient Space Pad:} Colchón atmosférico de texturas envolventes con retardo y reverberación expansiva.
    \item \textbf{Analog Bass 80s:} Bajo con pegada y riqueza sub-grave, emulando los sintetizadores de bajo de estado sólido.
    \item \textbf{Square Bell FM:} Campanas percusivas y cristalinas logradas con modulación de fase no entera.
    \item \textbf{Vocal AutoTune FX:} Ajustes optimizados para entrada de voz en vivo con corrección tonal asistida.
\end{enumerate}

\end{document}
